% Image credits for the photographs and AI illustrations of Book 4
% (University Chemistry, Year 3), French edition. Machine-readable license
% records live in images/book4/CREDITS.md; keep this file in sync with the
% English frontmatter/image-credits-book4.tex.
\chapter*{Crédits des images}

Les dessins de ce livre sont originaux. Les photographies sont utilisées
sous leurs licences libres respectives, avec nos remerciements à leurs
auteurs~; les liens vers les sources et les adresses des licences de
chaque photographie sont recensés dans
\texttt{images/book4/CREDITS.md}, dans le dépôt du livre.
Les pictogrammes de danger sont ceux du Système général harmonisé de
classification et d'étiquetage des produits chimiques, tels que les
dessine l'extension \texttt{ghsystem} de \TeX\ Live.

\bigskip
Outre les photographies et les dessins originaux, certaines figures sont
des illustrations générées par intelligence artificielle, produites avec
le générateur d'images d'OpenAI à partir de consignes rédigées par les
éditeurs du livre, et relues pour leur exactitude chimique avant
inclusion. La consigne complète de chaque image générée est consignée
dans \texttt{images/book4/ai/PROMPTS.md}, dans le dépôt.
\bigskip
\noindent Photographies, par chapitre~:
\begin{itemize}\small
  \item Chap.~1~: Erwin Schrödinger, 1933, Fondation Nobel, domaine public
        (Wikimedia Commons).
  \item Chap.~4~: cristaux de neige, Wilson A. Bentley (vers 1902), domaine
        public (Wikimedia Commons).
  \item Chap.~6~: antennes de radioastronomie sur un haut plateau,
        ESO/B.~Tafreshi (twanight.org), CC BY 4.0 (Wikimedia Commons).
  \item Chap.~7~: fluorine à la lumière du jour et sous lumière
        ultraviolette, par Masha Milshina, CC BY 4.0 (Wikimedia Commons).
  \item Chap.~8~: un laboratoire de RMN, par Shandchem, CC BY 2.0
        (Wikimedia Commons).
  \item Chap.~9~: William Lawrence Bragg, Fondation Nobel, domaine public
        (Wikimedia Commons).
  \item Chap.~10~: la tombe de Ludwig Boltzmann, par Feldkurat Katz,
        CC BY-SA 4.0 (Wikimedia Commons).
  \item Chap.~13~: spirales de Belousov--Zhabotinsky, par Stephen Morris,
        CC BY 2.0 (Wikimedia Commons).
  \item Chap.~14~: le trou de la couche d'ozone au-dessus de
        l'Antarctique, NASA, domaine public (Wikimedia Commons).
  \item Chap.~16~: micrographie électronique du nid d'abeilles d'un pot
        catalytique, par Galadrid, CC BY-SA 4.0 (Wikimedia Commons).
  \item Chap.~17~: rayons de soleil et effet Tyndall, par Kevin c higgins,
        CC BY-SA 4.0 (Wikimedia Commons).
  \item Chap.~18~: cristaux de rubis et d'émeraude, par Robert M. Lavinsky,
        CC BY-SA 3.0 (Wikimedia Commons).
  \item Chap.~20~: cristaux de ferrocène, par Andrea Sella, CC BY 4.0
        (Wikimedia Commons).
  \item Chap.~22~: un lingot de silicium monocristallin, par Sebastian
        Wallroth, domaine public (Wikimedia Commons).
  \item Chap.~23~: une fleur sur un aérogel de silice, NASA/JPL, domaine
        public~; la coupe de Lycurgue, par Marie-Lan Nguyen, CC BY 2.5
        (Wikimedia Commons).
  \item Chap.~24~: une limule de l'Atlantique, par Kaldari, CC0 (Wikimedia
        Commons).
  \item Chap.~27~: fleurs de pyrèthre, par Soramimi, CC BY-SA 4.0
        (Wikimedia Commons).
  \item Chap.~30~: écorce d'un if du Pacifique, par JOE BLOWE
        (theforestprimeval), CC BY-SA 2.0 (Wikimedia Commons).
  \item Chap.~32~: une prolifération d'algues vue par le satellite
        Landsat~8, USGS et NASA, domaine public.
  \item Chap.~33~: une double rampe à vide et à gaz inerte, par Polimerek,
        CC BY-SA 3.0 (Wikimedia Commons).
\end{itemize}
\clearpage
